\section{AegisKV Design}
\label{sec:design}

AegisKV separates complete recovery coverage from fast recovery access.
Shared DRAM and shared storage retain complete KV checkpoints; Peer HBM
holds selected hot blocks on a preassigned takeover DP. Its objective is
to reduce workflow completion time (WCT) under finite memory and I/O
capacity while preserving normal cache reuse.

\paragraph{Scope and assumptions.}
We consider fail-stop compute-worker failures with a compatible surviving
DP. The directory, ordered token log and advertised backing copies must
survive the protected failure domain; same-host DRAM or SSD does not
protect against loss of that host. The orchestrator exposes known
fork/join dependencies and accepted inputs. AegisKV preserves the engine's
native attention selections and replay semantics; it does not provide
exactly-once external tool side effects or protection against loss of all
advertised backing domains.

\subsection{Design Insights}
\label{sec:design-insights}

\paragraph{Value recovery in the workflow, without disabling normal reuse.}
A restore advances completion only when it shortens a required path
(Figure~\ref{fig:motivation-recovery-exposure}). However, a high-slack
agent may reuse its history immediately after a tool call. Assigning it
permanently to recomputation would consume shared serving capacity.
AegisKV therefore values both normal reuse and failure recovery, and
separates retaining any complete copy from paying for a faster copy.

\paragraph{Use sparse demand for acceleration, not completeness.}
Native sparse attention concentrates repeated historical KV references
on a subset of layer blocks. In Figure~\ref{fig:design-insights}, the
hottest approximately 20\% of blocks account for 87.64\%, 74.76\% and
56.54\% of accesses for NOSA-8B, NSA-8B and MiniCPM4.1-8B, respectively.
This concentration motivates selective Peer replicas, but does not
establish future-demand predictability. Later queries can change the
hot set, so a complete checkpoint and an on-demand miss path remain
necessary.

\begin{figure}[t]
 \centering
 \includegraphics[width=\columnwidth]{figures/kv-access-concentration.pdf}
 \caption{Native KV access concentration for NOSA-8B~\cite{huang2026nosav2},
 NSA-8B (MiniCPM-4-8B-NSA)~\cite{fasa2026minicpm4nsa}, and
 MiniCPM4.1-8B~\cite{openbmb2025minicpm41}. We use the final 512 queries
 of matched 32,704-token prefixes from three PG19 texts. Each query
 counts a historical 64-token block once across heads. Rank the 503
 blocks within each layer/window; average access shares over all
 32 layers, then texts. Bands show text min--max; the diagonal denotes
 uniform access. Compressed/selector state is excluded. This is
 retrospective concentration, not prediction accuracy or recovery speedup.}
 \label{fig:design-insights}
 \Description{Cumulative historical block access curves for NOSA-8B,
 NSA-8B and MiniCPM4.1-8B, averaged over every layer and three texts.
 About 20 percent of blocks account for 87.64, 74.76 and 56.54 percent
 of accesses, respectively; all three curves exceed the uniform baseline.}
\end{figure}

\paragraph{Pay only for transfer that remains exposed.}
Complete DRAM avoids historical reconstruction but still requires H2D.
A Peer copy can remove that transfer; later-layer transfers may already
be hidden by computation that must still execute. Useful residency
therefore depends on each request's committed history, uncommitted tail,
demand reveal time and recovery queues. The Master chooses per-request
layers and hot blocks; the Adapter controls only the total Peer budget.

\subsection{System Overview}
\label{sec:system-overview}

Figure~\ref{fig:system-overview} divides responsibilities across
Application, Controller and Engine. The Application supplies requests,
known dependencies and progress. The Controller contains
\textbf{(1)} Workflow Tracker, \textbf{(2)} Recovery-aware Router,
\textbf{(3)} KV Master and surviving metadata. Tracker supplies duration
estimates and slack; Router binds execution and takeover DPs; Master
chooses complete-checkpoint placement and Peer hot replicas. Within each
DP, \textbf{(4)} Adapter exposes a Peer budget and \textbf{(5)} Connector
executes transfers and recovery. Ordinary execution state, Peer copies
and transient workspace share physical HBM but have separate accounting.

\begin{figure}[t]
 \centering
 \includegraphics[width=\columnwidth]{figures/system-overview.pdf}
 \caption{AegisKV architecture. Black callouts match components (1)--(5)
 in the text; DP1 has the same Adapter/Connector as DP0. The example
 binds DP0 execution to DP1 takeover/Peer placement. Controller state
 includes manifests, token logs and leases. Blue links carry KV,
 dark dashes carry control, red dashes mark retry; orange local blocks
 illustrate an unsealed tail. Complete shared backing lies outside
 both DP failure boundaries. D/S denote DRAM/storage placement.}
 \label{fig:system-overview}
 \Description{Compact Application, Controller and Engine layers with
 numbered components, symmetric DP Local/Peer HBM and shared backing.}
\end{figure}

\paragraph{Objects and checkpoints.}
An immutable object identifies a sealed token block in one layer and TP
shard, together with model/weight version, layout, format and position
semantics. A legal group includes required selector state. Physical
blocks, sparse selection units and kernel tiles need not coincide; the
64-token unit above is a measurement convention, not a universal engine
requirement. Shared physical replicas consume capacity once.

A checkpoint manifest is complete only when every historical layer,
shard, auxiliary state and matching log prefix through frontier $M$ is
committed. A DRAM-resident manifest resolves all dependencies in DRAM;
it is not a hot subset split across tiers. Storage manifests are likewise
complete. Unchanged objects can be shared between checkpoints and only
new objects need copying, but partial extension cannot advance $M$.

Prefill offers newly sealed state; decode triggers review every configured,
block-aligned $Z$ tokens. Arrival, tool handoff, material demand/slack
changes and pressure also trigger review. The acknowledged tail $N$ can
exceed $Z$ when copying lags. Peer coverage is advertised only for objects
with complete committed backing in DRAM or storage.

\subsection{Workflow Tracker: Remaining Time and Slack}
\label{sec:workflow-protection}

Tracker turns application progress into versioned estimates of the
remaining workflow. A segment is an Agent's sequential continuation
between exposed fork/join boundaries or completion; it can contain
several LLM calls and tool waits. When an Agent waits for children, their
execution appears on child paths, not again in the parent's duration.
Completed work is removed. Unreleased segments retain precedence edges,
and independent tool waits remain explicit durations when exposed by
the application. Hidden future forks are not treated as known edges.

\paragraph{Estimate remaining execution, then derive slack.}
A GBDT maps a causal feature snapshot to seconds remaining in its segment.
Inputs are role/model, current input length, elapsed segment time,
completed-call count, and preceding calls' output lengths, API times,
context growth and gaps. Missing history is masked. Task/tool descriptors
are used only when supplied by the application. Training pairs a snapshot
with elapsed wall time from that snapshot to its recorded segment end;
training, calibration and evaluation partition whole workflows. Future
call counts, tool results and final output lengths are labels or future
events, never features. Elapsed time is not subtracted from a remaining-time
estimate again. Whole-Agent completion labels spanning nested forks are
not interchangeable with these segment labels.

A forward critical-path pass~\cite{kelley1959cpm} computes earliest
finishes; a backward pass from the predicted workflow finish computes
latest finishes. Their nonnegative difference is nominal segment slack.
Figure~\ref{fig:tracker-slack} gives a concrete example. Requests in one
segment share its margin; a new request does not receive an independent
slack allowance. Shared margins are handled by the Master's joint DAG
evaluation, rather than summed across requests.

\begin{figure}[t]
 \centering
 \includegraphics[width=\columnwidth]{figures/tracker-slack.pdf}
 \caption{From causal duration estimates to slack. In this illustrative
 residual DAG, B and C take 6 and 10\,s before a 2\,s continuation D.
 The predicted finish is 12\,s: nominal slack is 4\,s for B and zero for C.
 Refreshing B to 11\,s makes B critical and leaves C 1\,s of slack.
 Bold borders mark the initial critical path. These are constructed
 durations, not prediction-accuracy results.
 Publication of a new graph/model version invalidates old placement costs.}
 \label{fig:tracker-slack}
 \Description{Causal features feed a segment-duration predictor; a
 fork/join example derives slack and changes its critical branch on update.}
\end{figure}

\paragraph{Publish usable estimates and handle invalid ones.}
Each prediction carries a model version, feature-support mask, timestamp
and graph version. A calibrated model bundle also supplies a nonnegative
slack guard: an upper empirical quantile of positive slack overestimation
on its disjoint calibration workflows. Tracker subtracts this guard from
nominal slack and clips at zero. This is an empirical scheduling margin,
not a per-request probability guarantee. Master receives the original
durations and guard for joint evaluation; the clipped slack is an advisory
priority, not an additional delay to charge again. Model support, required edges,
calibration and maximum prediction age determine the validity flag.
If any check fails, Tracker publishes no positive slack credit; Master
uses the request-cost fallback in Section~\ref{sec:placement-value}.

Request arrival/completion, tool return and dependency changes refresh
the affected graph; periodic refresh bounds estimate age during long
calls. Each published update invalidates costs tied to the old version.
The two DAG passes take $O(|V|+|E|)$ in addition to predictor evaluation.
Duration errors affect placement quality, never checkpoint completeness
or replay semantics. Accuracy, false-positive slack and fallback frequency
are evaluated separately in Section~\ref{sec:eval-prediction}.

\subsection{Recovery-aware Router}
\label{sec:recovery-router}

Normal dispatch balances load with compatible prefix locality.
Figure~\ref{fig:agent-prefix-reuse} shows why Agent affinity is useful:
97.90\% potential same-agent token coverage and 0.37\% additional
cross-agent coverage in the trace. Cross-agent prefixes can still help
a subagent's first call, so affinity is a preference, not a restriction.
A tool handoff releases execution references and the request's logical
Peer pin; it does not delete reusable backing history.

\begin{figure}[t]
 \centering
 \includegraphics[width=\columnwidth]{figures/agent-prefix-reuse.pdf}
 \caption{Potential prefix reuse over 98,827 requests, 2,090 agents and
 393 CC/Weka workflows. Coverage uses earlier completed requests with
 the same recorded model and workflow. Same/other-agent coverage
 overlaps; only additional cross-agent coverage extends same-agent
 reuse. These are token-prefix opportunities, not runtime cache hits.}
 \label{fig:agent-prefix-reuse}
 \Description{Same-agent prefixes dominate aggregate reuse; cross-agent
 history can also help first subagent calls.}
\end{figure}

At arrival, Router records a versioned, layout-compatible takeover DP
outside the source failure domain. Peer copies reside on that target's
corresponding TP shards. On confirmed failure, Router fences the old
epoch, resolves a complete manifest and activates the bound target subject
to admission and workspace constraints. If it is unavailable, a new
binding can immediately use complete backing with zero Peer coverage;
bytes on the old target are not local hits on the new one.

Before output acknowledgement, a surviving ordered log records request
ID, token position/value and source epoch; prompt and accepted tool inputs
remain recoverable. Replay suppresses duplicate output and rejects stale
callbacks. Rebinding preserves the prior authoritative path until the new
binding is valid. Router queues and target load enter placement cost.

\subsection{KV Master: Three-Tier Placement and Direct Revaluation}
\label{sec:kv-master}

\subsubsection{Full-checkpoint placement and Peer admission}
\label{sec:placement-value}

The Master makes two coupled decisions: where to retain complete history,
and whether selected copies justify Peer capacity. At a checkpoint event,
new sealed state proposes a complete extension in DRAM, then storage if
DRAM cannot admit it. Admission may require demoting or evicting older
manifests. Unchanged immutable objects are shared; only absent physical
replicas require copying. If neither tier admits a complete extension,
$M$ remains unchanged and recovery replays a longer logged tail. A large
slack alone never suppresses an otherwise useful normal-reuse copy.

Candidate moves are complete storage-to-DRAM promotion, complete
DRAM-to-storage demotion, last-copy removal with reconstruction fallback,
and Peer hot-set replacement. The current placement is always an option.
Hard capacity relief is handled separately from optional acceleration.

For placement $P$ and a feasible candidate $P'$, Master evaluates
\begin{equation}
 \Delta J=\sum_e p_e\sum_w
       [T_w(P',e)-T_w(P,e)].
 \label{eq:value}
\end{equation}
Here $T_w(P,e)$ is the predicted remaining completion time of workflow
$w$. A finite scenario table describes normal continuation without a
fault, or the same continuation with a named source-domain fault at a
specified offset. Its probabilities $p_e$ sum to one. The table, decision
horizon and normal-reuse forecasts are deployment inputs, held fixed
across candidates; failure intensities are not substituted directly for
probabilities. Only admitted work and reuse events predicted from visible
history enter a forecast. A common fault includes all interrupted
requests, not independently sampled faults for each branch.

For each scenario, the cost model starts from the same queue and coverage
snapshot. It schedules proposed copies and requests with the Connector's
service order, measured service profiles and finite workspace. Complete
copy acknowledgements advance the modeled frontier; a fault offset fixes
the modeled logged tail. Recovery must wait for native layer dependencies
and any predicted demand misses. The resulting delays enter the residual
DAG as changes from the same ordinary-execution reference; normal work
already represented in the duration forecast is not added twice. Joins
take the maximum predecessor finish. Copy interference
and temporary occupancy affect those same completion times; no separate
byte-to-time price or second migration charge is added.

To apply the Tracker's slack guard jointly, a fault scenario adds the
guard once to each interrupted segment before the DAG pass, identically
for both placements. For one isolated branch this reduces its available
slack by the guard; for shared joins the full graph determines the effect.
The guard is never charged once per block or per request in that segment.
If the timing model is invalid, the same comparison instead sums predicted
normal-request and recovery-request completion delays from the resource
schedule, giving no credit for workflow slack. Invalid resource profiles
freeze optional promotion; they do not weaken safe demotion rules.

A common failure includes every interrupted request in that domain.
If two necessary branches are each 10\,s late, accelerating just one to
2\,s still leaves the join 10\,s late. Thus block scores cannot simply be
added. DRAM demotion is compared with actual storage-backed continuation;
Peer is compared with its complete lower path, including exposed H2D;
last-copy deletion is compared with an older checkpoint or reconstruction.
A backing demotion also changes the miss costs of dependent Peer plans.

\begin{figure}[t]
 \centering
 \includegraphics[width=\columnwidth]{figures/retention-value.pdf}
 \caption{Slack alone does not rank retention value. In this isolated-fault
 example, A uses two capacity units to reduce recovery from 4 to 2\,s;
 B uses one to reduce 18 to 8\,s. A at slack 0 saves 1\,s/unit, whereas
 B at slack 6 saves 10\,s/unit. Normal-use value is zero. These constructed
 profiles motivate cost/size comparison, not a policy-speedup claim.}
 \label{fig:retention-value}
 \Description{A higher-slack object can save more workflow time per byte.}
\end{figure}

\subsubsection{Choosing per-request layers and hot blocks}
\label{sec:sparse-placement}

For request $r$, $L_r$ is its Peer prefix depth and $K_{r,l}$ the number
of engine token blocks replicated in layer $l\le L_r$, including required
TP objects. The layer's hot set first includes native mandatory dependencies,
then blocks ranked by selection frequency over a configured window of
completed queries. Queries selecting one physical block through several
heads count once; ties favor recent access and then block ID. The engine's
mandatory set and selector state are preserved, and a budget below that
set's size is infeasible. Complete backing includes
the remaining blocks and selector state
(Figure~\ref{fig:sparse-recovery}). A prediction miss changes fetching,
not the engine's selection result.

\begin{figure}[t]
 \centering
 \includegraphics[width=\columnwidth]{figures/sparse-recovery.pdf}
 \caption{Complete backing, per-layer Peer budgets and demand misses.
 (a) The DRAM manifest covers B0--B3 in every layer; the tail is uncommitted.
 (b) On the bound target, $L=2$, with $K_1=2$ and $K_2=3$; layer 3 has
 no Peer copy. (c) Layer 1 needs B0/B2: B0 hits Peer and B2 is fetched
 from backing before computation. B3 was unused. The complete manifest
 also covers TP shards and selector state. If demoted as a whole,
 shared storage supplies the miss path before H2D.}
 \label{fig:sparse-recovery}
 \Description{Three complete historical layers, an uncommitted tail,
 different hot-block counts and an explicitly routed layer-1 miss.}
\end{figure}

The candidate generator uses a finite list of budget templates. A template
specifies a block count or fraction for each layer; thus $K$ can remain
fixed or grow with history without changing the algorithm. For each
template, Master considers prefix depths from zero to model depth,
discarding infeasible mandatory sets. This costs model-depth times
template-count candidates, rather than a Cartesian product of all layer
budgets. The current placement is included even if it is off-grid.
At a join with tied bottleneck predecessors, it also constructs bounded
bundles that advance all tied branches to their next deeper candidate
under one template. Bundles exceeding the configured request limit are
skipped; every attempt counts toward the common work budget.

Each candidate uses the current $M$; future $N$ is derived separately in
each fault-offset scenario, including commit lag beyond $Z$. Causal
selection statistics supply preload and miss demand; an unsupported
sparse profile uses full historical demand. Only computation that must
still execute hides transfer. This bounded search is a heuristic over
explicit complete plans; it makes no global-optimality claim.

Figure~\ref{fig:design-pipeline} illustrates why layer order, prediction
and Peer residency complement each other. The matched comparison is
(c) versus (d): both retain complete DRAM backing and replay the same
64-token tail. Peer saves 24\,ms of exposed waiting in this model; it
neither removes miss fetches nor avoids the required replay computation.
The first seven layers reach the model's demand-only wait floor at this
fixed profile, not a universal depth or a bound under contention.

\begin{figure}[t]
 \centering
 \includegraphics[width=\columnwidth]{figures/recovery-pipeline.pdf}
 \caption{Analytical layered recovery with and without Peer.
 MiniCPM4-8B has 32 layers, 32K history and a
 64-token tail. Replay is calibrated at 2.870\,ms/layer; assumed effective
 H2D bandwidth is 6.25\,GB/s (5.369\,ms per full layer). Trace-derived
 mean preload and preload-plus-miss fractions are 0.6584/0.7748.
 (a) waits for bulk restoration; (b) overlaps full-layer restoration;
 (c) preloads hot blocks and fetches revealed misses; (d) adds 147.5\,MiB
 of Peer hot sets in the first seven layers. Numbers are tail completion
 and exposed wait. Cold copies follow replay; the ideal preemptible
 stream and homogeneous layer costs omit bounded-workspace and queue
 effects. This is not end-to-end fault-recovery measurement.}
 \label{fig:design-pipeline}
 \Description{The same required replay computation with progressively
 reduced exposed KV waits: 172, 83, 44 and 20 milliseconds.}
\end{figure}

Neither $L$ nor $K$ must increase on every decode step. A fixed sparse
budget can preserve $K$ while changing block identities; a proportional
budget can increase demand, but Peer admission remains capacity-limited.
At each $Z$ review or material event, the Master can retain or replace the
current set. Proposed and committed sets remain distinct until transfers
finish. If no exposed wait is removed, $L_r=0$ is valid even at zero slack.

\subsubsection{Online admission, direct revaluation, and eviction}
\label{sec:placement-plan}

Each decision snapshots DAG and binding versions, causal demand, complete
manifests, queues, leases and physical byte ledgers. DRAM/storage candidates
are complete manifests; Peer candidates are legal layer/block bundles.
The safe victim set excludes physical readers and copy dependencies.
Normal Peer competition also respects running requests' logical pins;
request completion or tool handoff unpins them. Self hot-set replacement
and Adapter pressure may override logical pins, never physical leases.

For each candidate, a bounded rotating scan examines eligible victims
in stable-ID order. It appends the feasible removal with least current
whole-plan loss per physical byte released, breaking ties by fewer moved
bytes and then ID. The tentative placement is revalued after every
removal. A manifest releasing no bytes alone is skipped unless a legal
owner bundle can release its shared objects. A DRAM victim needing
storage space includes lower-tier relief and fallback changes in the
same plan. Search stops when space exists for the candidate and its
workspace, or the work budget is exhausted.

The ratio constructs a victim set; the final decision evaluates the
complete candidate-plus-victim plan with Equation~\ref{eq:value}.
Optional placement requires $\Delta J<-\epsilon$, where $\epsilon$ is a
configured time margin that suppresses churn; migration itself is already
in the resource schedule. Required capacity relief has no incoming
candidate: among examined feasible plans, it chooses the smallest cost
increase, even if positive. If no safe plan fits, admission pauses;
an infeasible plan never changes coverage or makes reserved bytes free.

\par\medskip\noindent\begin{minipage}{\columnwidth}
\raggedright
 \refstepcounter{designalgorithm}\label{alg:placement}
 \textbf{Algorithm \thedesignalgorithm: Direct placement revaluation}
 \par\smallskip
 \begin{tabularx}{\columnwidth}{@{}r@{\quad}X@{}}
 \toprule
 1 & $S\gets\textsc{Snapshot}();\quad P^*\gets\varnothing;\quad d^*\gets\infty$ \\
 2 & $h\gets\textsc{ShrinkRequired}(S)$ \\
 3 & \textbf{for} $m\in\textsc{Candidates}(S,h)$ \textbf{while} \textsc{WorkLeft}() \\
 4 & \quad $P\gets\textsc{MakeRoom}(S,m)$ \\
 5 & \quad \textbf{if} $P=\varnothing$ \textbf{then continue} \\
 6 & \quad $d\gets\Delta J(P;S)$ \\
 7 & \quad \textbf{if} $d<d^*$ \textbf{then} $(P^*,d^*)\gets(P,d)$ \\
 8 & \textbf{if} $P^*=\varnothing$ \textbf{then return} \\
 9 & \textbf{if} $\neg h\land d^*\ge-\epsilon$ \textbf{then return} \\
 10 & \textbf{if not} \textsc{ReserveIfCurrent}$(S,P^*)$ \textbf{then} \\
 11 & \quad \textsc{Replan}(); \textbf{return} \\
 12 & \textsc{Submit}$(P^*)$ \\
 \bottomrule
 \end{tabularx}
\end{minipage}\par\medskip

In Algorithm~\ref{alg:placement}, $h$ selects capacity relief rather than
optional placement. \textsc{MakeRoom} uses the victim scan above and returns
a feasible complete plan, or $\varnothing$ if none fits within the work
budget. \textsc{ReserveIfCurrent} rechecks versions and leases before
reserving executable capacity; \textsc{Submit} invokes the Connector
protocol. Completion or failure reconciles state and triggers replanning.
An early return retains state and pauses admission if capacity is insufficient.

Cached estimates are reused only while their dependencies remain valid.
Profiling, victim comparisons and DAG passes all count against the work
budget. Unsupported or stale resource profiles freeze optional promotions;
necessary relief uses bounded recency order under the same safety rules.
Direct revaluation carries no rent, credit or renewal state. Evaluation
compares its decision and control costs with a rental alternative
(Section~\ref{sec:eval-rental}).

\subsection{Adapter: Windowed Peer Capacity}
\label{sec:adapter}

Each DP's Adapter sets a Peer budget for the next time window from the
managed HBM pool, expected serving allocation and reserved headroom,
capped by a configured maximum and clipped at zero. Serving allocation
includes execution state and near-term demand, without subtracting Peer
occupancy twice. The window duration and cap are configurable; $Z$ is
fixed independently. Adapter does not set per-request $L$ or $K$.

A lower target pauses new Peer admission while actual plus reserved
occupancy exceeds the target. Master selects a safe relief plan, including
logically pinned replicas when necessary. A valid immutable backing copy
allows clean Peer release without redundant D2H. Active reads or copies
must drain before memory becomes available. Emergency foreground pressure
can request relief within a window; if leases prevent immediate release,
serving applies backpressure rather than reusing occupied addresses.

\subsection{Connector: Committed Copies and Layered Recovery}
\label{sec:connector}

Connector applies reservations and bounded transfer queues for foreground
reuse, demanded recovery and checkpoint movement. Configured service
shares and transfer quanta limit interference; real misses precede
speculative reads within the recovery share. Copying and Peer occupancy
remain costs in the placement objective.

\paragraph{Publish complete state before releasing its predecessor.}
A full extension commits only after validation of every required layer,
TP shard, selector object and corresponding token-log prefix. Moving a
complete DRAM checkpoint to storage follows copy, backend persistence,
manifest commit, reader redirection, reader drainage and source release.
Partial writes cannot release the source; duplicate or stale-epoch
callbacks cannot advance coverage. Temporary duplicate objects and
workspace count against real capacity, and a buffered write acknowledgement
alone is not proof of the backend's persistence contract.

Before last-copy deletion, Connector publishes a complete older manifest
or a prompt-and-token-log reconstruction fallback. Peer claims lacking
complete backing are withdrawn in the same transition. On copy failure,
the prior authoritative location remains valid and failed reservations
are reconciled; pending copies never count as current protection.

\paragraph{Restore according to actual demand.}
Recovery acquires the manifest, binding and location leases. Compatible
Peer blocks on the bound DP bypass H2D; missing blocks use complete DRAM,
or shared-storage reads followed by H2D. Layer-ordered prefetch overlaps
only with remaining replay computation. Actual selection misses block
the affected operation until inputs arrive; dense layers use full demand
and missing selector state is restored before selection. TP collectives
and replay batching preserve native semantics. Subsequent decode uses the
same demand-miss path for historical blocks still absent from HBM.
Adopted Peer objects move
to execution accounting without double-counting or freeing their bytes.
Prefill failures use the last complete frontier, possibly zero.

\paragraph{A recovery across the three tiers.}
Consider DP0 failing while a new checkpoint is being copied. Router fences
DP0's epoch and selects its bound DP1. The incomplete new copy is ignored;
Connector acquires the last committed manifest and leases its locations.
If that checkpoint is in DRAM, missing Peer blocks use H2D. If its whole
manifest has already moved to storage, they use storage reads followed by
H2D. Committed Peer objects still bypass those transfers. Connector
replays the acknowledged log beyond $M$, suppresses duplicate output,
then resumes decode. If DP1 cannot admit the execution working set, the
request queues or rebinds; a new target starts with no assumed Peer hits.

\paragraph{Why prediction errors do not corrupt recovery.}
Three invariants separate placement from correctness. First, every
advertised checkpoint resolves complete, version-compatible dependencies
and its log prefix. Second, a location switch publishes the complete new
manifest before retiring its predecessor; removing the last copy first
publishes a logged reconstruction path. Third, physical memory is released
only after dependent transfers and readers drain. The initial log-only
state satisfies these conditions. Creating a pending copy does not change
coverage; a validated commit extends it; redirection and drainage preserve
it during demotion or deletion. A failed copy leaves the predecessor in
force. Thus these transitions preserve the invariants under the stated
surviving-directory/log assumptions, regardless of Slack or hot-set error.
Such errors can add waiting, copying or recomputation; they cannot authorize
missing KV as a hit. Progress additionally requires a compatible target
and eventually available execution/I/O capacity.
